% Generated draft: requires mathematical and source review.
\documentclass[11pt]{article}
\usepackage[margin=25mm]{geometry}
\usepackage{fontspec}
\setmainfont{Latin Modern Roman}
\usepackage{amsmath,amssymb,mathtools,unicode-math}
\setmathfont{Latin Modern Math}
\usepackage{enumitem,needspace,xurl,hyperref}
\setlength{\parindent}{0pt}
\setlength{\parskip}{5pt}
\setlength{\emergencystretch}{3em}
\newcommand{\N}{\mathbb{N}}
\newcommand{\Z}{\mathbb{Z}}
\newcommand{\Q}{\mathbb{Q}}
\newcommand{\R}{\mathbb{R}}
\newcommand{\C}{\mathbb{C}}
\newcommand{\F}{\mathbb{F}}
\newcommand{\A}{\mathbb{A}}
\newcommand{\PP}{\mathbb{P}}
\newcommand{\E}{\mathbb{E}}
\newcommand{\eps}{\varepsilon}
\newcommand{\dd}{\,\mathrm{d}}
\begin{document}
\section*{UnsolvedMath ID: 30002515}
\textbf{Source problem code:} catalogue entry
\subsection{Definitions and mathematical statement}
For an integer genus \(g\geq 0\), let \(S_g\) be a closed, connected, oriented surface of genus \(g\), and let \(\operatorname{Mod}_g\) denote its group of isotopy classes of orientation-preserving diffeomorphisms. The action on integral first homology gives the symplectic representation
\[
\sigma_g:\operatorname{Mod}_g\longrightarrow\operatorname{Sp}(2g,\Z),
\]
where a symplectic basis identifies \(H_1(S_g;\Z)\) with \(\Z^{2g}\), and \(\operatorname{Sp}(2g,\Z)\) preserves its intersection pairing. Write \(\mathcal M_g\) for the moduli object of smooth compact connected complex curves of genus \(g\); its precise interpretation and the intended range of \(g\) are not specified in the source.

For a holomorphic map \(\varphi:X\to\mathcal M_g\) with the mapping-class monodromy assumed in the source, write
\[
\rho_\varphi:\pi_1(X,x_0)\longrightarrow\operatorname{Mod}_g,
\qquad
\rho'_\varphi=\sigma_g\circ\rho_\varphi.
\]
A base point and identifications of reference homology groups are needed to express these as particular homomorphisms; without such choices, monodromy is ordinarily specified only up to conjugacy.

\textbf{Source question.} For every finite-type Riemann surface \(X\) and every pair of holomorphic maps \(\varphi_1,\varphi_2:X\to\mathcal M_g\) having the indicated monodromy, does
\[
\rho'_{\varphi_1}=\rho'_{\varphi_2}
\quad\Longrightarrow\quad
\varphi_1(x)=\varphi_2(x)\ \text{for every }x\in X
\]
hold?

\textbf{Unresolved qualifications.} The source does not specify the genus range, the precise meaning of finite type, or whether \(\mathcal M_g\) is a coarse moduli space, an orbifold, or a stack. In particular, it does not explain the additional data needed to obtain mapping-class monodromy from a map to a coarse moduli space. It also leaves unresolved whether equality of monodromy means literal equality after compatible choices or equality up to conjugacy, and what equality of maps means if a stack is intended. No nonconstancy or nonisotriviality hypothesis is supplied; none is imposed here. These issues require clarification before this question can be regarded as a fully specified rigidity assertion.
\subsection{Short English statement}
The source asks whether a holomorphic map from a finite-type Riemann surface to curve moduli is determined by its monodromy action on integral first homology. The genus range, moduli interpretation, comparison of monodromy, and possible restrictions on the maps remain unspecified. The supplied open label is dated curation, not an independently verified mathematical conclusion.
\subsection{Sources}
\begin{itemize}
\item Local UnsolvedMath registry: supplied complete record for permanent ID \(30002515\), titled \emph{Determination of Moduli Maps by Symplectic Monodromy}, in the Oberwolfach Reports Open Problems collection. Exact source problem code: \texttt{catalogue entry}. The record gives proposed year 2014 and does not name an individual proposer.
\item Supplied primary provenance: \emph{New Trends in Teichmüller Theory and Mapping Class Groups} (2014), Oberwolfach Reports, pp.~395--452, DOI \texttt{10.4171/owr/2014/07}. Listed in the supplied literature sources as \emph{Oberwolfach Report 07/2014}: \url{https://doi.org/10.4171/owr/2014/07}. The supplied background describes these reports as documenting mathematical workshops and their open-problem discussions; no exact problem page is supplied.
\item Additional supplied literature source: \emph{Holomorphic maps between configuration spaces}, \url{https://doi.org/10.1007/s00209-025-03893-6}.
\item Literature assessment dated \texttt{2026-08-21}: the supplied curation labels the problem open and states that rigidity results are known when full mapping-class monodromy agrees, whereas determination after passing only to the symplectic representation remains unresolved in the literature checked. This is supplied dated curation derived from the OWR material, not independently verified here; the background explicitly requests independent verification of time-sensitive status and cited claims.
\item The registry labels the statement \texttt{corrected\_verified} and describes its changes as normalization of notation, typography, and wording without changing mathematical content. This supplied editorial label does not resolve the formulation issues identified above or constitute mathematical review of this draft.
\end{itemize}
\end{document}
